\centerline{\bf \Huge Геометрия I}
\centerline{\bf \Huge Вписанные углы}

\begin{flushright}
        \scriptsize{Выживает тот, у кого есть воля к жизни.\\Мисато Кацураги. Евангелион 24 серия.}
\end{flushright}

\problem{1}
Окружности $S_1$ и $S_2$ пересекаются в точках $A$ и $P$. Через точку $A$ проведена касательная $AB$ к окружности $S_11$, а через точку $P$ --- прямая $CD$, параллельная $AB$ (точки $B$ и $C$ лежат на $S_2$, точка $D$ --- на $S_1$). Докажите, что $ABCD$ --- параллелограмм.

\problem{2}
Окружности $S_1$ и $S_2$ пересекаются в точке $A$. Через точку $A$ проведена прямая, пересекающая $S_1$ в точке $B$, а $S_2$ в точке $C$. В точках C и B проведены касательные к окружностям, пересекающиеся в точке $D$. Докажите, что угол $BDC$ не зависит от выбора прямой, проходящей через $A$.

\problem{3}
Вписанная окружность касается сторон $AB$ и $AC$ треугольника $ABC$ в точках $M$ и $N$. Пусть $P$ --- точка пересечения прямой MN и биссектрисы угла B (или её продолжения). Докажите, что $\angle BPC = 90^{\circ}$

\problem{4}
В окружность вписаны треугольники $T_1$ и $T_2$, причём вершины треугольника $T_2$ являются серединами дуг, на которые окружность разбивается вершинами треугольника $T_1$. Докажите, что в шестиугольнике, являющемся пересечением треугольников $T_1$ и $T_2$, диагонали, соединяющие противоположные вершины, параллельны сторонам треугольника $T_1$ и пересекаются в одной точке.


\problem{5}\textit{Лемма о бабочке.}
Через середину $C$ произвольной хорды $AB$ окружности проведены две хорды $KL$ и MN (точки $K$ и $M$ лежат по одну сторону от $AB$). Отрезки $KN$ и $ML$ пересекают $AB$ в точках $Q$ и $P$. Докажите, что $PC=QC$.